% ============================================================
% EJERCICIO 7 — Comparación de herramientas de IA (Módulo 2)
% Pida la MISMA tabla de 4x4 (resistencias) a dos o más IA con el MISMO prompt.
% Datos para la tabla (use estos):
%   Resistor | Nominal (ohm) | Medido (ohm) | Error (%)
%   R1 | 100 | 98,5 | 1,5 ;  R2 | 220 | 221,3 | 0,6
%   R3 | 330 | 327,9 | 0,6 ;  R4 | 470 | 475,2 | 1,1
% ============================================================
\documentclass[12pt,a4paper]{article}
\usepackage[utf8]{inputenc}
\usepackage[T1]{fontenc}
\usepackage[spanish]{babel}
\usepackage{amsmath,amssymb}
\usepackage{booktabs}
\usepackage{siunitx}
\sisetup{locale=FR, output-decimal-marker={,}}
\usepackage{circuitikz}
\usepackage[hidelinks]{hyperref}
\usepackage{fancyvrb}

\begin{document}
\section{Prompt utilizado}
\begin{Verbatim}[frame=single,fontsize=\small]
TODO el mismo prompt para todas las herramientas
\end{Verbatim}

\section{Resultados por herramienta}
% Repita la subsección por herramienta.
\subsection{TODO Herramienta A}
\begin{Verbatim}[frame=single,fontsize=\small]
TODO código generado (sin editar)
\end{Verbatim}
\subsection{TODO Herramienta B}
\begin{Verbatim}[frame=single,fontsize=\small]
TODO código generado (sin editar)
\end{Verbatim}

\section{Comparación}
\begin{table}[htbp]
  \centering
  \caption{Comparación de herramientas}
  \label{tab:comparacion}
  \begin{tabular}{lccc}
    \toprule
    Criterio & Herramienta A & Herramienta B & Herramienta C \\
    \midrule
    Compila sin cambios (sí/no) & TODO & TODO & TODO \\
    Usa \texttt{booktabs}/\texttt{siunitx} & TODO & TODO & TODO \\
    Errores a corregir (n.°) & TODO & TODO & TODO \\
    Caption y label & TODO & TODO & TODO \\
    Puntaje (1-5) & TODO & TODO & TODO \\
    \bottomrule
  \end{tabular}
\end{table}

\section{Reflexión}
% TODO ¿Cuál funcionó mejor y por qué? ¿Qué cambiaría en el prompt?
\end{document}
% RÚBRICA (20 pts): mismo prompt y código sin editar (5) · tabla comparativa con
% criterios objetivos (6) · capturas adjuntas (3) · reflexión fundamentada (6)
